\par\indent\SourceLeaf{205}%
Dieses Resultat lässt sich auch so ausdrücken: Indem die Charakteristik \(f\) eine beliebige Function bedeutet, hat man den reellen Theil von \(f(x+iy)\) für \(X\), und den imaginären Theil, mit Weglassung des Factors \(i\), entweder für \(Y\) oder für \(-Y\) anzunehmen.

Gebraucht man die Charakteristiken \(\varphi,\varphi'\) in der Bedeutung des Art. 7 und setzt
\[
\varphi(x+iy)=\xi+i\eta,\quad\varphi'(x-iy)=\xi-i\eta
\]
\noindent wo offenbar \(\xi\) und \(\eta\) reelle Functionen von \(x\) und \(y\) sein werden, so hat man, in der \textit{ersten} Auflösung,
\[
\begin{aligned}
\mathrm{d}X+i\,\mathrm{d}Y&=(\xi+i\eta)(\mathrm{d}x+i\,\mathrm{d}y)\\
\mathrm{d}X-i\,\mathrm{d}Y&=(\xi-i\eta)(\mathrm{d}x-i\,\mathrm{d}y)
\end{aligned}
\]
\noindent und folglich
\[
\begin{aligned}
\mathrm{d}X&=\xi\,\mathrm{d}x-\eta\,\mathrm{d}y\\
\mathrm{d}Y&=\eta\,\mathrm{d}x+\xi\,\mathrm{d}y
\end{aligned}
\]
\noindent Macht man nun
\[
\begin{aligned}
\xi&=\sigma\mathord{.}\cos\gamma,&\eta&=\sigma\mathord{.}\sin\gamma\\
\mathrm{d}x&=\mathrm{d}s\mathord{.}\cos g,&\mathrm{d}y&=\mathrm{d}s\mathord{.}\sin g\\
\mathrm{d}X&=\mathrm{d}S\mathord{.}\cos G,&\mathrm{d}Y&=\mathrm{d}S\mathord{.}\sin G
\end{aligned}
\]
\noindent so dass \(\mathrm{d}s\) ein Linearelement in der ersten Ebne, \(g\) dessen Neigung gegen die Abscissenlinie, \(\mathrm{d}S\) das correspondirende Linearelement in der zweiten Ebne und \(G\) dessen Neigung gegen die Abscissenlinie bedeutet, so geben die obigen Gleichungen
\[
\begin{aligned}
\mathrm{d}S\mathord{.}\cos G&=\sigma\mathord{.}\mathrm{d}s\mathord{.}\cos(g+\gamma)\\
\mathrm{d}S\mathord{.}\sin G&=\sigma\mathord{.}\mathrm{d}s\mathord{.}\sin(g+\gamma)
\end{aligned}
\]
\noindent und folglich, wenn man. was erlaubt ist, \(\sigma\) als positiv betrachtet,
\[
\mathrm{d}S=\sigma\mathord{.}\mathrm{d}s,\quad G=g+\gamma
\]
\noindent Man sieht also (in Uebereinstimmung mit Art. 7), dass \(\sigma\) das Verhältniss der Vergrösserung des Elements \(\mathrm{d}s\) in der Darstellung \(\mathrm{d}S\) vorstellt, und. wie gehörig, von \(g\) unabhängig ist; und eben so zeigt die Unabhängigkeit des Winkels \(\gamma\) von \(g\), dass alle von einem Punkte ausgehende Linearelemente in der ersten Ebne
